%% reference: etex
% The errors a scanner raises when the value it reads is *wrong* rather than
% absent. Each is a distinct `tex.web` section with a distinct recovery, which
% is why the sequence is the test: getting a message right but its recovery
% wrong shows up as the next error moving, or as an error that should not be
% there at all.
%
% Judged against eTeX, because one of them is about register range and that is
% the single place this engine diverges from TeX82 on purpose — 32768 registers
% rather than 256, because expl3 allocates far past 256. Against `tex -ini`,
% `\count300` would be an error here and is not; against eTeX both agree.
\catcode`\{=1 \catcode`\}=2

% §459: a dimension whose number has no unit. TeX assumes `pt` and *keeps the
% number*, so this is 5pt, not 0pt. We already computed 5pt — silently. The
% value was right and the silence was wrong.
\dimen0=5
\message{PROBE:nounit=\the\dimen0}

% §433 as widened by eTeX to `scan_register_num`: out of range is an error and
% the recovery is *zero*, not a wrap. We wrapped modulo 32768, so this used to
% set \count7232 and leave \count0 alone.
\count40000=1
\message{PROBE:badreg=\the\count0}

% ...and inside the range, no error: this is the divergence from TeX82.
\count300=7
\message{PROBE:etexreg=\the\count300}

% §434 `scan_char_num`: 0..255, recovery zero. We kept 300.
\chardef\c=300
\message{PROBE:badchar=\the\c}

% §436 `scan_fifteen_bit_int`: a mathchar packs class, family and position, so
% its range is wider — 300 is perfectly ordinary here.
\mathchardef\mc=300
\message{PROBE:mathchar=\the\mc}

% §434/§1227: the target of every character-code table is still an eight-bit
% TeX character, and an invalid target recovers to slot zero. The right side
% then has a table-specific range and recovers to value zero. `\delcode` alone
% accepts negative values.
\lccode256=65
\message{PROBE:badcodetarget=\the\lccode0}
\lccode65=256
\message{PROBE:badlccode=\the\lccode65}
\catcode1=16
\message{PROBE:badcatcode=\the\catcode1}
\mathcode66=32769
\message{PROBE:badmathcode=\the\mathcode66}
\sfcode67=-1
\message{PROBE:badsfcode=\the\sfcode67}
\delcode68=16777216
\message{PROBE:baddelcode=\the\delcode68}
\delcode69=-123
\message{PROBE:negdelcode=\the\delcode69}

% §436 is a separate restricted scanner with its own diagnostic and help.
\mathchardef\badmc=32768
\message{PROBE:badmathchar=\the\badmc}

% §454: third-order infinity is the ceiling, and every extra `l` is its own
% error. Recovery keeps the component at `filll`.
\skip2=1pt plus 2filllll
\message{PROBE:filll=\the\skip2}

% §459 again inside a rule specification, where the recovery has to leave the
% rule with the number it read rather than a default.
\setbox0\hbox{\vrule width 3}
\message{PROBE:rule=\the\wd0}

% §435 `scan_four_bit_int`, reached from §1350: an output stream is 0..15 and
% the recovery is zero. Only `\openout` gets this check — `\closeout` is scanned
% with `\write`'s clamp instead, because §1350 branches on node size and
% `close_node` happens to be the same size as `write_node`. So the second line
% below raises nothing, and its silence is as much the test as the first line's
% error: getting it wrong in the safe-looking direction adds an error TeX does
% not have.
\immediate\openout-1=zzz
\immediate\closeout99
\message{PROBE:stream=ok}

% and the neighbouring error, for the boundary between them: no number at all
% is §446, not §459.
\count1=x
\message{PROBE:letter=\the\count1}
\end
